\documentclass[10pt,a4paper]{amsart}
\usepackage[margin=19mm]{geometry}
\usepackage{amsmath,amssymb,mathtools}
\usepackage[scr=rsfs,cal=euler]{mathalfa}
\usepackage{xfrac}
\usepackage[unicode,hidelinks,bookmarks=false]{hyperref}
\newcommand{\ie}{i.e.\ }
\newcommand{\cf}{cf.\ }
\newcommand{\resp}{respectively\ }
\newcommand{\Subsection}{\subsection}
\providecommand{\nobreakdash}{\nobreak\hbox{-}}
\setlength{\emergencystretch}{2em}
\multlinegap0pt
\usepackage{bm}
\usepackage{bbm}
\usepackage{dsfont}
\usepackage{xspace}
\usepackage{cancel}
\usepackage{mathtools}
\usepackage{amsthm}
\makeatletter
\@ifundefined{theorem}{\newtheorem{theorem}{Theorem}}{}
\@ifundefined{lemma}{\newtheorem{lemma}[theorem]{\bf Lemma}}{}
\makeatother
\DeclareMathOperator{\supp}{supp}
\newcommand{\Abs}[1]{\left\vert #1 \right\vert}
\newcommand{\R}{\mathbb{R}}
\newcommand{\abs}[1]{\vert #1 \vert}
\newcommand{\norm}[1]{\left\Vert #1 \right\Vert}
\title[A Regularized Shallow Water System]{A Regularized Shallow Water System}
\author{Evgueni Dinvay, Henrik Kalisch}
\date{Source: arXiv:2605.30044v1 (CC BY 4.0)}
\begin{document}
\maketitle

\noindent\emph{Source and licence.} A Regularized Shallow Water System. Version arXiv:2605.30044v1 (CC BY 4.0), distributed under the Creative Commons Attribution 4.0 International licence, as recorded in the arXiv licence metadata for this exact version and shown on the abstract page https://arxiv.org/abs/2605.30044v1. Full rights notice: licenses/arxiv-2605.30044-CC-BY-4.0.txt.\par\smallskip

\noindent\textit{Editorial scope.} Counted unit: the Lemma whose source label is \texttt{None} in the article (source0.tex lines 859--870, source label \texttt{None}), delivered together with the article's own complete original proof of it (source0.tex lines 872--1626, one \texttt{proof} environment of 755 lines). No mathematics is written, repaired or re-derived: statement and proof are the article's own text line for line. Required context retained uncounted: regularizedSW (lines 259--273); Hamiltonian\_structure (lines 311--324); variational\_derivatives (lines 326--342); u\_is\_Au (lines 482--485); local\_well\_posedness\_lemma (lines 492--505); A\_and\_f\_definition (lines 662--690); energy\_norm (lines 818--824). Every definition, display and label of those lines is retained unchanged. Omitted from this package and not redistributed: 1--258, 274--310, 325--325, 343--481, 486--491, 506--661, 691--817, 825--858, 871--871, 1627--2514. Nothing inside a retained excerpt is omitted or altered: every retained body is delivered exactly as the article's lines are written, so no omission receipt of any kind accompanies this package. Citations: the retained excerpts carry no citation command at all, so no bibliography entry of the article is transcribed and no reference is redistributed. Reference adaptation: none was needed and none was made; every reference inside the delivered unit resolves inside it, so no unresolved reference or citation is produced. Printed slips kept literal: no printed word, symbol, hypothesis or number of the counted statement or of its proof was altered. Layout and class: the article's own class is \texttt{article} with packages including amscd, amssymb, verbatim, tikz, graphicx, a4wide, amsfonts, latexsym, amsthm, amsmath, overpic, natbib, url, hyperref; the wrapper is the lane's pinned \texttt{amsart} editorial wrapper at 10pt on A4, which restates the theorem-like environments the retained text needs and adds the article's own macro definitions that the retained text needs (\texttt{\textbackslash Abs}, \texttt{\textbackslash R}, \texttt{\textbackslash abs}, \texttt{\textbackslash norm}, \texttt{\textbackslash supp}), with extra wrapper packages (bm, bbm, dsfont, xspace, cancel, mathtools). The delivered numbering is the unit's own. Assets: no retained span contains a figure environment, a graphics-inclusion command, a PDF-page inclusion, a table or a TikZ picture; no image file is copied into this package, the package declares no asset dependency and no image file is redistributed with it. Licence: version arXiv:2605.30044v1 of this article is distributed under the Creative Commons Attribution 4.0 International licence, as recorded in the arXiv licence metadata for this exact version and shown on the abstract page https://arxiv.org/abs/2605.30044v1. A full rights notice is delivered at licenses/arxiv-2605.30044-CC-BY-4.0.txt. Characters: every retained excerpt is pure ASCII, so the delivered engine, XeLaTeX with its default Latin Modern text font, reports no missing character.
\begin{equation}
\label{regularizedSW}
	\left.
	\begin{array}{rcl}
		\eta_t + h_0 u_x  + K ( \eta u )_x
		& = &
		0
		, \\
		u_t + g \eta_x + K ( u^2/2 )_x
		& = &
		0
		,
	\end{array}
	\right\}
\end{equation}

\begin{equation}
\label{Hamiltonian_structure}
    \eta_t
    =
    - K \partial_x
    \frac{\delta \mathcal H}{\delta u}
    \quad
    \text{ and }
    \quad
    u_t
    =
    - K \partial_x
    \frac{\delta \mathcal H}{\delta \eta}
\end{equation}

\begin{equation}
\label{variational_derivatives}
    \frac{\delta \mathcal H}{\delta \eta}
    =
    g K^{-1} \eta
	+
    \frac 12 u^2
    \quad
    \text{ and }
    \quad
    \frac{\delta \mathcal H}{\delta u}
    =
    h_0 K^{-1} u
    +
    \eta u
    .
\end{equation}

\begin{equation}
\label{u_is_Au}
	 v = \mathcal A(v; v_0)
\end{equation}

\begin{lemma}[local well-posedness]
\label{local_well_posedness_lemma}
	Let $s > 1/2$, $v_0 = (\eta_0, u_0)^T \in X^s$
	and $T = ( 7C_s \lVert v_0 \rVert_{X^s} )^{-1}$
	with some constant $C_s > 0$ depending only on $s$.
	Then there exists a unique solution
	$v = (\eta, u)^T \in X^s_T$ of Problem \eqref{u_is_Au}.


	Moreover, for any $R > 0$ there exists $T = T(R) > 0$
	such that the flow map associated with Equation \eqref{u_is_Au}
	is a real analytic mapping of the open ball
	$B_R(0) \subset X^s$ to $X^s_T$.
\end{lemma}

\begin{equation}
\label{A_and_f_definition}
    A
    =
    \begin{pmatrix}
        0
        &
        -\partial_x
        \\
        -\partial_x
        &
        0
    \end{pmatrix}
    \quad
    \text{ and }
    \quad
    f(v)
    =
    f(\eta, u)
    =
	-
    K \partial_x
	\begin{pmatrix}
		\eta u
		\\
		u^2 / 2
	\end{pmatrix}
    .
\end{equation}

\begin{equation}
\label{energy_norm}
    \norm{v}_{\mathcal H}
    =
    \frac 1{\sqrt 2}
    \norm{K^{-1/2} v}_{L^2}
\end{equation}

\begin{lemma}
The energy functional $\mathcal H$ is continuously differentiable on $X^{1/2}$.
Moreover,
\[
    \mathcal H(v(t))
    =
    \mathcal H(v_0)
    , \quad
    t \in [0, T],
\]
where $v, v_0$ and $T$ are described in Lemma \ref{local_well_posedness_lemma}.
\end{lemma}

\begin{proof}
%
Linearisation of $\mathcal H$ around $v \in X^{1/2}$
gives
\[
    d \mathcal H(v) \delta v
    =
	\int_\R
	\left(
		\eta K^{-1} \delta \eta
		+
        u K^{-1} \delta u
		+
        \frac 12 \delta \eta u^2
        +
        \eta u \delta u
	\right)
	dx
    , \quad
    \delta v = ( \delta \eta, \delta u )^T \in X^{1/2}
    ,
\]
which defines a bounded linear operator $d \mathcal H(v)$,
due to
\begin{multline*}
    \Abs{d \mathcal H(v) \delta v}
    \leqslant
    \norm{
		K^{-1/2} \eta
    }_{L^2}
    \norm{
        K^{-1/2} \delta \eta
    }_{L^2}
    +
    \norm{
        K^{-1/2} u
    }_{L^2}
    \norm{
        K^{-1/2} \delta u
    }_{L^2}
    \\
    +
    \frac 12
    \norm{
        \delta \eta
    }_{L^2}
    \norm{
        u
    }_{L^4}^2
    +
    \norm{
        \eta
    }_{L^4}
    \norm{
        u
    }_{L^4}
    \norm{
        \delta u
    }_{L^2}
    \lesssim
    \left(
        1 + \norm{v}_{X^{1/2}}^2
    \right)
    \norm{\delta v}_{X^{1/2}}
    ,
\end{multline*}
and so $\mathcal H$ is differentiable.
Its gradient with respect to the energy inner product, inducing the norm \eqref{energy_norm},
has the form
\[
    \nabla \mathcal H(v)
    =
    2 K
    \frac{ \delta \mathcal H }{ \delta v }
    =
    \left(
        2 K
        \frac{ \delta \mathcal H }{ \delta \eta }
        ,
        2 K
        \frac{ \delta \mathcal H }{ \delta u }
    \right)^T
    =
    \begin{pmatrix}
        2 \eta + K \left( u^2 \right)
        \\
        2 u + 2 K \left( \eta u \right)
    \end{pmatrix}
    ,
\]
where the variational derivatives are defined in \eqref{variational_derivatives}.
Continuity of the derivative $d \mathcal H$ is equivalent to
the continuity of the  gradient $\nabla \mathcal H$.
The latter follows from the estimate
\begin{multline*}
    \norm{ \nabla \mathcal H(v_1) - \nabla \mathcal H(v_2) }_{\mathcal H}
    \leqslant
    2
    \norm{ v_1 - v_2 }_{\mathcal H}
    +
    \frac 1{\sqrt 2}
    \norm{
        \begin{pmatrix}
            u_1^2 - u_2^2
            \\
            2 \eta_1 u_1 - 2 \eta_2 u_2
        \end{pmatrix}
    }_{L^2}
    \\
    \lesssim
    \left(
        1
        +
        \norm{ v_1 }_{\mathcal H}
        +
        \norm{ v_2 }_{\mathcal H}
    \right)
    \norm{ v_1 - v_2 }_{\mathcal H}
    ,
\end{multline*}
where we have used the boundedness of $K^{1/2}$ in $L^2(\R)$,
H\"older's inequality
\(
    \norm{fg}_{L^2}
    \leqslant
    \norm{f}_{L^4}
    \norm{g}_{L^4}
\)
and the Sobolev embedding $H^{1/4}(\R) \hookrightarrow L^4(\R)$.
This completes the proof of smoothness of $\mathcal H$.


In order to prove the energy conservation,
we introduce a regularized version  of System \eqref{regularizedSW} as follows.
%
%
For $\mu \geqslant 1$ let $P_\mu$
be the Fourier multiplier with an even smooth symbol $\theta_\mu(\xi) = \theta(\xi/\mu)$,
that is, for $f \in \mathcal S'(\R)$,
\[
    \widehat{P_\mu f}(\xi) = \theta\left( \frac{\xi}{\mu} \right) \widehat f(\xi)
    .
\]
We assume that $\supp\theta \subset [-2,2]$,
hence $\widehat{P_\mu f}$ is supported in $[-2\mu,2\mu]$, so $P_\mu f(x)$ is smooth.
We also assume that $\abs{\theta} \leqslant 1$.
Then $P_\mu$ is bounded in $H^s(\R)$ for any $s$, with operator norm $\leqslant 1$. 
Finally, we assume that $\theta = 1$ on $[-1, 1]$,
so that $P_\mu$ converges strongly to the identity operator
in $H^s(\R)$ as $\mu \to \infty$.


Now consider the following frequency-truncated version
\begin{equation}
\label{frequency_truncatedSW}
    \partial_t v = A P_{\mu} v + f(v)
    , \quad
    v(0) = v_0
\end{equation}
of the Cauchy problem for System \eqref{regularizedSW}
with the initial data $v_0$,
where $A$ and $f$ are defined in \eqref{A_and_f_definition}.
As above we can introduce the unitary group
\(
    \mathcal S_{\mu}(t) = e^{A P_{\mu} t}
\)
and notice that a solution of \eqref{frequency_truncatedSW}
satisfies the integral equation
%
\begin{equation}
\label{frequency_truncated_duhamel}
	v(t)
    =
    \mathcal S_{\mu}(t - t_0) v(t_0)
	+
    \int_{t_0}^t \mathcal S_{\mu}(t - r) f(v(r))
	dr
    , \quad
    0 \leqslant t_0 \leqslant t
	.
\end{equation}
%
On $X^{1/2}$ we define
%
\begin{equation*}
%\label{Hamiltonian}
	\mathcal H_{\mu}(v)
    =
	\mathcal H_{\mu}(\eta, u)
    =
    \frac 12 \int_\R
	\left(	
		\eta K^{-1} P_{\mu} \eta
		+
        u K^{-1} P_{\mu} u
		+
        \eta u^2
	\right)
	dx
\end{equation*}
%
that is obviously smooth.
It is straightforward to repeat the above arguments and conclude
that there exists a unique solution
\[
    v = v_{\mu}
    \in
    C   \left( [0, T]; X^s \right)
    \cap
    C^1 \left( (0, T); X^s \right)
\]
of \eqref{frequency_truncatedSW} for every given $\mu$.
Importantly,
here $s$ and $T$ are exactly as in the formulation of Lemma \ref{local_well_posedness_lemma}.
In particular, $T$ is independent of $\mu$.
Clearly,
\(
	\mathcal H_{\mu}(v(t))
\)
is smooth with respect to $t \in (0, T)$ and the chain rule
\[
    \frac d{dt}
	\mathcal H_{\mu}(v(t))
    =
    d \mathcal H_{\mu}(v(t))
    \partial_t v(t)
    =
	\int_\R
	\left(
		\eta K^{-1} P_{\mu} \partial_t \eta
		+
        u K^{-1} P_{\mu} \partial_t u
		+
        \frac 12 u^2 \partial_t \eta
        +
        \eta u \partial_t u
	\right)
	dx
\]
applies.
It is easy to see that \eqref{frequency_truncatedSW} has
the Hamiltonian structure \eqref{Hamiltonian_structure}
with the new functional $\mathcal H_{\mu}$ in place of $\mathcal H$
and the new variational derivatives
\begin{equation*}
    \frac{\delta \mathcal H_{\mu}}{\delta \eta}
    =
    K^{-1} P_{\mu} \eta
	+
    \frac 12 u^2
    \quad
    \text{ and }
    \quad
    \frac{\delta \mathcal H_{\mu}}{\delta u}
    =
    K^{-1} P_{\mu} u
    +
    \eta u
    .
\end{equation*}
%
Therefore,
\begin{multline*}
    \frac d{dt}
	\mathcal H_{\mu}(v(t))
    =
	\int_\R
	\left(
        \frac{\delta \mathcal H_{\mu}}{\delta \eta}
        \partial_t \eta
		+
        \frac{\delta \mathcal H_{\mu}}{\delta u}
        \partial_t u
	\right)
	dx
    =
    -
	\int_\R
	\left(
        \frac{\delta \mathcal H_{\mu}}{\delta \eta}
        K \partial_x
        \frac{\delta \mathcal H_{\mu}}{\delta u}
		+
        \frac{\delta \mathcal H_{\mu}}{\delta u}
        K \partial_x
        \frac{\delta \mathcal H_{\mu}}{\delta \eta}
	\right)
	dx
    \\
    =
    -
	\int_\R
    \partial_x
	\left(
        \frac{\delta \mathcal H_{\mu}}{\delta \eta}
        K
        \frac{\delta \mathcal H_{\mu}}{\delta u}
	\right)
	dx
    =
    0
    ,
\end{multline*}
%
and so $\mathcal H_{\mu}(v(t))$ is constant.
By continuity,
\(
    \mathcal H_{\mu}(v(t)) = \mathcal H_{\mu}(v_0)
\)
for all $t \in [0, T]$,
where
$v = v_{\mu}$ is the solution of the regularized Cauchy problem \eqref{frequency_truncatedSW}.


Now let $v$ be the solution of the original Cauchy problem described
by Lemma \ref{local_well_posedness_lemma}
and $v_{\mu}$ be the solution of \eqref{frequency_truncatedSW}.
We will show that $v_{\mu} \to v$ in $L^2(0, T; X^s)$ as $\mu \to \infty$.
The idea is to split the whole interval $[0, T]$ into small sub-intervals
of size $h > 0$ and to prove this convergence repeatedly on every sub-interval
moving from left to right.
Both $v$ and $v_{\mu}$ satisfy the integral equation \eqref{frequency_truncated_duhamel}
with $\mathcal S$ and $\mathcal S_{\mu}$ in place of the semigroup, respectively.
Hence
%
\begin{equation*}
	v_{\mu}(t)
    -
	v(t)
    =
    \mathcal S_{\mu}(t - t_0) v_{\mu}(t_0)
    -
    \mathcal S(t - t_0) v(t_0)
	+
    \int_{t_0}^t
    \left(
        \mathcal S_{\mu}(t - r) f(v_{\mu}(r))
        -
        \mathcal S(t - r) f(v(r))
	\right)
    dr
\end{equation*}
%
for $0 \leqslant t_0 \leqslant t \leqslant t_0 + h \leqslant T$
with $h$ to be chosen below independently of $\mu$.
We estimate this difference as
%
\begin{multline*}
	\norm{
        v_{\mu}(t)
        -
    	v(t)
    }_{X^s}
    \leqslant
	\norm{
        v_{\mu}(t_0)
        -
    	v(t_0)
    }_{X^s}
    +
	\norm{
        \left(
            \mathcal S_{\mu}(t - t_0)
            -
            \mathcal S(t - t_0)
        \right)
        v(t_0)
    }_{X^s}
	\\
    +
    \int_{t_0}^t
    \norm{
        f(v_{\mu}(r))
        -
        f(v(r))
	}_{X^s}
    dr
    +
    \int_{t_0}^t
    \norm{
        \left(
            \mathcal S_{\mu}(t - r)
            -
            \mathcal S(t - r)
    	\right)
         f(v(r))
	}_{X^s}
    dr
    .
\end{multline*}
%
Squaring this inequality and then integrating the result from $t_0$ to $t_0 + h$, one obtains
%
\begin{multline*}
    \frac 1{16}
	\norm{
        v_{\mu}
        -
    	v
    }_{ L^2(t_0, t_0 + h; X^s) }^2
    \leqslant
    h
	\norm{
        v_{\mu}(t_0)
        -
    	v(t_0)
    }_{X^s}^2
    \\
    +
    \int_{t_0}^{t_0 + h}
	\norm{
        \left(
            \mathcal S_{\mu}(t - t_0)
            -
            \mathcal S(t - t_0)
        \right)
        v(t_0)
    }_{X^s}^2
    dt
    +
    h^2
    \int_{t_0}^{t_0 + h}
    \norm{
        f(v_{\mu}(r))
        -
        f(v(r))
	}_{X^s}^2
    dr
	\\
    +
    h
    \int_{t_0}^{t_0 + h}
    \int_{t_0}^{t_0 + h}
    \norm{
        \left(
            \mathcal S_{\mu}(t - r)
            -
            \mathcal S(t - r)
    	\right)
         f(v(r))
	}_{X^s}^2
    dr dt
    .
\end{multline*}
%
From the proof of Lemma \ref{local_well_posedness_lemma} one can easily deduce
\(
	\norm{
        v_{\mu}
    }_{X_T^s}
    ,
	\norm{
    	v
    }_{X_T^s}
    \leqslant
	3
    \norm{
        v_0
    }_{X^s}
    ,
\)
and so
\[
    h^2
    \int_{t_0}^{t_0 + h}
    \norm{
        f(v_{\mu}(r))
        -
        f(v(r))
	}_{X^s}^2
    dr
    \leqslant
    C h^2
	\norm{
        v_{\mu}
        -
    	v
    }_{ L^2(t_0, t_0 + h; X^s) }^2
    ,
\]
where $C$ is independent of $\mu$.
This implies
%
\begin{multline*}
    \frac 1{20}
	\norm{
        v_{\mu}
        -
    	v
    }_{ L^2(t_0, t_0 + h; X^s) }^2
    \leqslant
    h
	\norm{
        v_{\mu}(t_0)
        -
    	v(t_0)
    }_{X^s}^2
    \\
    +
    \int_{t_0}^{t_0 + h}
    \int_{\mathbb R}
    \Abs{
        \mathcal F
        \left[
            \left(
                \mathcal S_{\mu}(t - t_0)
                -
                \mathcal S(t - t_0)
            \right)
            v(t_0)
        \right]
        (\xi)
	}^2
    \left( 1 + \xi^2 \right)^s
    d \xi
    dt
	\\
    +
    h
    \int_{t_0}^{t_0 + h}
    \int_{t_0}^{t_0 + h}
    \int_{\mathbb R}
    \Abs{
        \mathcal F
        \left[
            \left(
                \mathcal S_{\mu}(t - r)
                -
                \mathcal S(t - r)
        	\right)
             f(v(r))
        \right]
        (\xi)
	}^2
    \left( 1 + \xi^2 \right)^s
    d \xi
    dr dt
    .
\end{multline*}
%
These two integrals tend to zero as $\mu \to \infty$ by the dominated convergence theorem,
since
\[
    e^{ \pm i (t - r) \xi \theta_{\mu}(\xi) }
    -
    e^{ \pm i (t - r) \xi }
    \to
    0
    \text{ as }
    \mu \to \infty
    \text{ for every }
    t, r, \xi
\]
and the integrands are bounded by
\[
    2
    \Abs{
        \mathcal F
        \left[
            v(t_0)
        \right]
        (\xi)
	}^2
    \left( 1 + \xi^2 \right)^s
    \quad
    \text{ and }
    \quad
    2
    \Abs{
        \mathcal F
        \left[
             f(v(r))
        \right]
        (\xi)
	}^2
    \left( 1 + \xi^2 \right)^s
    ,
\]
%
respectively.
The first term
\(
	\norm{
        v_{\mu}(t_0)
        -
    	v(t_0)
    }_{X^s}
    \to
    0
\)
by choosing $t_0$ properly at least for a subsequence $\mu = \mu_k \to \infty$.
Indeed, for the first subinterval $t_0 = 0$ and this term is zero.
Therefore,
\(
    v_{\mu} \to v
\)
in
\(
    L^2(0, h; X^s)
\)
and so a subsequence
\(
    \norm{v_{\mu_k} - v}_{X^s}
    \to
    0
\)
almost everywhere on $[0, h]$.
Next candidate $t_0$ for left boundary of subinterval is taken in $(h/2 ,h]$,
so that
\(
    \norm{v_{\mu_k}(t_0) - v(t_0)}_{X^s}
    \to
    0
    .
\)
Thus by induction we obtain
\(
    \norm{v_{\mu_k} - v}_{X^s}
    \to
    0
\)
in $L^2(0, T)$
and almost everywhere on $[0, T]$.


Finally, we can prove the energy conservation
by analyzing the difference
%
\begin{equation}
\label{energy_conservation_difference}
    \mathcal H(v(t))
    -
    \mathcal H(v_0)
    =
    \mathcal H(v(t))
    -
    \mathcal H_{\mu}(v_{\mu}(t))
    +
    \mathcal H_{\mu}(v_0)
    -
    \mathcal H(v_0)
\end{equation}
%
for $t \in [0, T]$ and with $\mu = \mu_k$ described above.
Clearly,
\[
    \Abs{
        \mathcal H_{\mu}(v_0)
        -
        \mathcal H(v_0)
    }
    \leqslant
    \norm{v_0}_{\mathcal H}
    \norm{( P_{\mu} - I )v_0}_{\mathcal H}
    \to
    0
    .
\]
The first difference in \eqref{energy_conservation_difference}
is bounded by
\[
    \Abs{
        \mathcal H(v(t))
        -
        \mathcal H_{\mu}(v_{\mu}(t))
    }
    \leqslant
    \Abs{
        \mathcal H(v(t))
        -
        \mathcal H_{\mu}(v(t))
    }
    +
    \Abs{
        \mathcal H_{\mu}(v(t))
        -
        \mathcal H_{\mu}(v_{\mu}(t))
    }
    ,
\]
where the first term is estimated in the same way
\[
    \Abs{
        \mathcal H(v(t))
        -
        \mathcal H_{\mu}(v(t))
    }
    \lesssim
    \norm{( P_{\mu} - I )v(t)}_{\mathcal H}
    \to
    0
    ,
\]
while the last term is estimated by the mean-value theorem as follows.
The modified energy $\mathcal H_{\mu}$ is a continuously differentiable functional,
which one can check by repeating the arguments given
for analysis of $\mathcal H$ at the beginning of the proof.
Moreover,
\[
    \norm{d \mathcal H_{\mu}(w)}
    \leqslant
    C
    \left(
        1 + \norm{w}_{X^{1/2}}^2
    \right)
    \leqslant
    C
    \left(
        1 + \norm{w}_{X^s}^2
    \right)
\]
%
with the constant $C$ independent of $\mu$.
Substituting $w = \lambda v(t) + (1 - \lambda) v_{\mu}(t)$
with $0 < \lambda < 1$
and recalling that both $v(t)$ and $v_{\mu}(t)$ are bounded independently of $\mu$,
one obtains
\[
    \Abs{
        \mathcal H_{\mu}(v(t))
        -
        \mathcal H_{\mu}(v_{\mu}(t))
    }
    \leqslant
    C
    \left(
        1 + 9 \norm{v_0}_{X^s}^2
    \right)
    \norm{
        v(t)
        -
        v_{\mu}(t)
    }_{X^s}
    .
\]
%
The latter converges for almost every $t \in [0, T]$.
Thus, passing to the limit $\mu \to \infty$ in \eqref{energy_conservation_difference},
we obtain
\(
    \mathcal H(v(t))
    -
    \mathcal H(v_0)
    =
    0
\)
for almost every $t \in [0, T]$.
The proof is concluded
by continuity of
\(
    \mathcal H(v(t))
    .
\)
%
\end{proof}

\end{document}
